\section*{Bab \ref{ch:g10:lewis-and-shape} --- Struktur Lewis dan Bentuk Molekul}

\begin{solution}{exo:g10:lewis-and-shape:1}
Sepasang \omterm{def:g9:inside-the-atom:nucleus}{elektron} yang digunakan bersama oleh dua \omterm{def:g7:atoms-and-molecules:atom}{atom}. Sepasang
\omterm{def:g10:electron-shells:valence}{elektron valensi} yang hanya dimiliki satu \omterm{def:g7:atoms-and-molecules:atom}{atom}, tidak digunakan bersama.
\end{solution}

\begin{solution}{exo:g10:lewis-and-shape:2}
H: 1 (kurang satu \omterm{def:g9:inside-the-atom:nucleus}{elektron} untuk duplet). C: 4, N: 3, O: 2 (kurang 4, 3,
dan 2 \omterm{def:g9:inside-the-atom:nucleus}{elektron} untuk oktet).
\end{solution}

\begin{solution}{exo:g10:lewis-and-shape:3}
\chemfig{H-\charge{0=\:,90=\:,270=\:}{Cl}}: klorin memiliki pasangan bersama dan tiga
\omterm{def:g10:lewis-and-shape:lone-pair}{pasangan bebas}, 8 \omterm{def:g9:inside-the-atom:nucleus}{elektron}.
\end{solution}

\begin{solution}{exo:g10:lewis-and-shape:4}
$5 + 3 \times 1 = 8$ \omterm{def:g9:inside-the-atom:nucleus}{elektron}, 4 pasang (tiga ikatan, satu \omterm{def:g10:lewis-and-shape:lone-pair}{pasangan bebas}).
\end{solution}

\begin{solution}{exo:g10:lewis-and-shape:5}
\ce{CH4}: tetrahedron. \ce{NH3}: piramida. \ce{H2O}: bengkok.
\ce{CO2}: linear.
\end{solution}

\begin{solution}{exo:g10:lewis-and-shape:6}
\chemfig{H-\charge{90=\:,270=\:}{S}-H}: seperti air, empat kelompok di sekeliling belerang,
dua di antaranya \omterm{def:g10:lewis-and-shape:lone-pair}{pasangan bebas}: bengkok.
\end{solution}

\begin{solution}{exo:g10:lewis-and-shape:7}
$2 \times 5 = 10$ \omterm{def:g9:inside-the-atom:nucleus}{elektron}, 5 pasang. Satu ikatan tunggal dan 4
\omterm{def:g10:lewis-and-shape:lone-pair}{pasangan bebas} membuat setiap nitrogen hanya memiliki 6 \omterm{def:g9:inside-the-atom:nucleus}{elektron};
mengubah dua \omterm{def:g10:lewis-and-shape:lone-pair}{pasangan bebas} menjadi pasangan bersama menghasilkan
\omterm{def:g10:lewis-and-shape:multiple-bond}{ikatan rangkap tiga} dan satu \omterm{def:g10:lewis-and-shape:lone-pair}{pasangan bebas} pada setiap \omterm{def:g7:atoms-and-molecules:atom}{atom}:
\chemfig{\charge{180=\:}{N}~\charge{0=\:}{N}}.
\end{solution}

\begin{solution}{exo:g10:lewis-and-shape:8}
\chemfig{H-C(-[2]H)(-[6]H)-\charge{90=\:,270=\:}{O}-H}: empat ikatan di sekeliling karbon,
tetrahedral; di sekeliling oksigen dua ikatan dan dua \omterm{def:g10:lewis-and-shape:lone-pair}{pasangan bebas},
bengkok.
\end{solution}

\begin{solution}{exo:g10:lewis-and-shape:9}
Pada \ce{CO2}, karbon memiliki dua kelompok (dua \omterm{def:g10:lewis-and-shape:multiple-bond}{ikatan rangkap dua}) dan
tidak ada \omterm{def:g10:lewis-and-shape:lone-pair}{pasangan bebas}: keduanya mengarah berlawanan. Pada \ce{H2O},
oksigen memiliki empat kelompok (dua ikatan, dua \omterm{def:g10:lewis-and-shape:lone-pair}{pasangan bebas}) dalam
susunan tetrahedral: kedua ikatannya membentuk sudut.
\end{solution}

\begin{solution}{exo:g10:lewis-and-shape:10}
\chemfig{C(-[3]H)(-[5]H)=C(-[1]H)-[7]H}: di sekeliling setiap karbon,
tiga kelompok (satu \omterm{def:g10:lewis-and-shape:multiple-bond}{ikatan rangkap dua}, dua ikatan tunggal): segitiga
datar, dengan sudut mendekati \ang{120}.
\end{solution}

\begin{solution}{exo:g10:lewis-and-shape:11}
Karbon di pusat, empat ikatan tunggal ke \omterm{def:g7:atoms-and-molecules:atom}{atom-atom} klorin yang
masing-masing memiliki tiga \omterm{def:g10:lewis-and-shape:lone-pair}{pasangan bebas}: empat kelompok di sekeliling
karbon, sebuah tetrahedron.
\end{solution}

\begin{solution}{exo:g10:lewis-and-shape:12}
Etanol, \chemfig{H-C(-[2]H)(-[6]H)-C(-[2]H)(-[6]H)-\charge{90=\:,270=\:}{O}-H}, dan
metoksimetana,
\chemfig{H-C(-[2]H)(-[6]H)-\charge{90=\:,270=\:}{O}-C(-[2]H)(-[6]H)-H}.
\end{solution}

\begin{solution}{exo:g10:lewis-and-shape:13}
\omterm{def:g10:lewis-and-shape:lone-pair}{Pasangan bebas} memerlukan ruang sedikit lebih besar daripada \omterm{def:g10:lewis-and-shape:lone-pair}{pasangan
ikatan} dan mendorong ikatan-ikatan saling mendekat. Amonia memiliki satu
\omterm{def:g10:lewis-and-shape:lone-pair}{pasangan bebas}, air dua: sudutnya menyempit dari metana ke amonia ke air.
\end{solution}

\begin{solution}{exo:g10:lewis-and-shape:14}
\chemfig{H-C~\charge{0=\:}{N}}: dua kelompok di sekeliling karbon, linear.
\chemfig{H-C(-[6]H)=\charge{45=\:,315=\:}{O}}: tiga kelompok di sekeliling karbon, segitiga
datar.
\end{solution}

\begin{solution}{exo:g10:lewis-and-shape:15}
$3 \times 6 = 18$ \omterm{def:g9:inside-the-atom:nucleus}{elektron}, 9 pasang:
\chemfig{\charge{180=\:,270=\:}{O}=\charge{90=\:}{O}-\charge{0=\:,90=\:,270=\:}{O}}. Oksigen yang di tengah
memiliki tiga kelompok (satu \omterm{def:g10:lewis-and-shape:multiple-bond}{ikatan rangkap dua}, satu ikatan tunggal,
satu \omterm{def:g10:lewis-and-shape:lone-pair}{pasangan bebas}): \omterm{def:g7:atoms-and-molecules:molecule}{molekulnya} bengkok.
\end{solution}

\begin{solution}{pb:g10:lewis-and-shape:1}
\textbf{1.} \chemfig{H-C(-[2]H)(-[6]H)-H}, \chemfig{H-\charge{90=\:}{N}(-[6]H)-H},
\chemfig{H-\charge{90=\:,270=\:}{O}-H}.

\textbf{2.} Empat kelompok pada masing-masing.

\textbf{3.} Tidak ada pada metana, satu pada amonia, dua pada air.

\textbf{4.} Karbon membentuk empat ikatan; oksigen membentuk dua (dua
kelompok lainnya adalah \omterm{def:g10:lewis-and-shape:lone-pair}{pasangan bebas}, yang tidak digambarkan dengan
tongkat dalam perangkat model).

\textbf{5.} \chemfig{H-C(-[2]H)(-[6]H)-C(-[2]H)(-[6]H)-H}: satu tongkat
menghubungkan kedua bola karbon, dan diperlukan enam bola hidrogen.

\textbf{6.} Tetrahedron, piramida, bengkok.

\textbf{7.} \omterm{def:g10:lewis-and-shape:lone-pair}{Pasangan bebas} memerlukan lebih banyak ruang daripada
\omterm{def:g10:lewis-and-shape:lone-pair}{pasangan ikatan} dan menekan ikatan-ikatan saling mendekat; air, dengan
dua \omterm{def:g10:lewis-and-shape:lone-pair}{pasangan bebas}, memiliki sudut terkecil.

\textbf{8.} Lubang-lubang bola karbon membentuk sudut \ang{109.5}; sudut
air yang sebenarnya \ang{104.5}: bola oksigen memerlukan dua lubangnya
sendiri pada \ang{104.5}.

\textbf{9.} \ang{180}: dua kelompok di sekeliling karbon mengarah
berlawanan, dan \omterm{def:g7:atoms-and-molecules:molecule}{molekulnya} linear.

\textbf{10.} Dua titik sudut yang tidak pada rusuk yang sama adalah titik
sudut yang berhadapan pada satu sisi: jaraknya adalah diagonal sisi,
$\sqrt{a^2 + a^2} = a\sqrt{2}$.

\textbf{11.} Pusat kubus terletak di tengah-tengah diagonal ruang:
$a\sqrt{3}/2$.

\textbf{12.} H--H: $\sqrt{2} \approx \qty{1.41}{cm}$; C--H:
$\sqrt{3}/2 \approx \qty{0.87}{cm}$.

\textbf{13.} Setiap hidrogen berada di sebuah titik sudut kubus dan
karbon di pusatnya: setiap titik sudut berjarak sama dari pusat.

\textbf{14.} Segitiga C\,H$_1$\,H$_2$ sama kaki (C\,H$_1$ = C\,H$_2$);
garis dari puncaknya ke titik tengah alasnya tegak lurus pada alas itu.

\textbf{15.} $\sin \widehat{\text{M\,C\,H}_1} = \dfrac{\text{M\,H}_1}{\text{C\,H}_1}
= \dfrac{a\sqrt{2}/2}{a\sqrt{3}/2} = \sqrt{2/3} \approx 0.816$.

\textbf{16.} Setengah sudut $\approx \ang{54.7}$; sudut H--C--H $\approx
2 \times 54.7 = \ang{109.5}$.

\textbf{17.} Itulah sudut \ang{109.5} tetrahedron pada tabel.

\textbf{18.} Lubang-lubang itu harus dibor dengan sudut \ang{109.5} satu
sama lain.
\end{solution}
